\section{Introduction}

Generative AI is reshaping decision systems by changing not only how information is produced, but also how responsibility, oversight, and authority are distributed between humans and machines. In many contemporary domains, the critical organizational question is no longer whether AI should be used at all, but what role it should occupy in the decision pipeline. An AI system may provide advice, recommend a decision for human confirmation, act as a default that humans can override, or operate with substantial autonomy. These are not merely interface choices. They are institutional choices about delegation, supervision, accountability, and governance.

Existing scholarship has generated important insights into human responses to AI and automation. Research on trust in automation emphasizes that the key challenge is not blind acceptance but appropriate reliance \citep{lee2004trust}. Work on algorithm aversion shows that people may penalize algorithms heavily after witnessing errors \citep{dietvorst2015algorithm}, whereas work on algorithm appreciation suggests that under some task conditions people prefer algorithmic over human judgment \citep{logg2019algorithm}. A related body of research on AI-assisted decision making investigates how confidence scores, explanations, and performance cues calibrate trust and improve local decision quality \citep{zhang2020confidence,ma2023trust}. Yet these literatures largely analyze pointwise reactions: whether a system is trusted, whether advice is adopted, or whether case-level reliance is calibrated.

This paper argues that the more consequential problem lies one level higher. Humans do not only learn whether AI is reliable. They also learn what role AI should occupy in a decision system. The outcome of this learning process is not simply a revised belief about accuracy, but a revised collaboration architecture. Over repeated feedback, people may converge toward stable patterns of low or high AI authorization. Crucially, these trajectories may be path dependent: early AI errors may trigger stronger negative updating than comparable human errors, reducing later willingness to expose the system to AI-led configurations and thereby slowing recovery even when long-run AI performance equals that of human decision makers.

The central claim of this research prospectus is therefore straightforward: the key consequence of AI errors is not merely lower short-run trust, but altered convergence paths in human--AI collaboration architecture learning. To formalize this argument, the paper adapts the logic of experience-weighted attraction learning (EWA) \citep{camerer1999ewa}. Here, the strategic objects are not game-theoretic moves in the conventional sense, but collaboration architectures that differ in autonomy, supervision, cost, and risk exposure. This move allows us to connect repeated feedback, asymmetric learning from AI versus human errors, path dependence, architecture convergence, and behavioral steady states within a single dynamic framework.

The project contributes to ongoing debates in human--AI collaboration by shifting the analytical lens from adoption to role formation, from single advice episodes to long-run architecture convergence, and from learning about AI capability to learning about AI role. The empirical design proposed below is built to detect precisely these dynamics.
